\subsection{10. Ein korrigierter, wirklich entscheidbarer
Anschluss-Test}\label{ein-korrigierter-wirklich-entscheidbarer-anschluss-test}

Ein sinnvoller nächster Versuch besteht nicht nur aus einem
unbeschrifteten \codeword{P U\_AB P}. Er sollte zusammen liefern:

\begin{enumerate}
\def\labelenumi{\arabic{enumi}.}
\tightlist
\item
  \textbf{Globale Quelle und Einbettungen:} konkrete primitive Algebra
  und zwei CAR-/Tensor-erhaltende Abbildungen; alle Kreuzrelationen und
  die tatsächliche globale Graduierung.
\item
  \textbf{Gemeinsamer Zustand:} derselbe globale H und Zustand;
  Nachweis, welche bisherigen lokalen Pole und Lücken darin noch gelten.
  Kein stiller Rückgriff auf den Produktgrundzustand unabhängiger
  Banken.
\item
  \textbf{Physische Operation:} aus einem erklärten Quellenwort
  erzeugter aktiver Operator oder Generator, mit Zeit- beziehungsweise
  Ressourcenskala. Passive Chartwechsel bleiben getrennt.
\item
  \textbf{Gram-bereinigte Dynamik:} \codeword{S\_ij=<h\_i|h\_j>} und
  \codeword{K\_ij=<h\_i|H|h\_j>} berechnen; das generalisierte
  Eigenproblem beziehungsweise eine orthonormalisierte Darstellung
  verwenden. Prüfen, ob mehr als \codeword{K=E\_h S} entsteht.
\item
  \textbf{Operativer Transfer:} dieselbe Anfangspräparation und
  Zielmessung; zeitabhängige Änderung gegenüber der statischen
  Überlappung und Fehlerkontrolle gegenüber dem übrigen Zustandsraum.
  Erst dann mit dem bekannten bedingten Zwei-Banken-Transfer
  vergleichen.
\end{enumerate}

Das nimmt die mögliche gemeinsame Herkunft ernst, korrigiert aber die
falschen Automatismen. Der stärkste sofortige Erkenntnisgewinn dieses
Audits ist negativ und präzise: \textbf{Allein durch Überlappung oder
Umbenennung eines entarteten nativen Polraums entsteht keine
Transportdynamik.} Ob der Compiler darüber hinaus genau die passende
globale Dynamik besitzt, ist die konkrete offene Frage.

\begin{center}\rule{0.5\linewidth}{0.5pt}\end{center}

\section{Anhang E: arithmetischer
Schleifenvorschlag}\label{anhang-e-arithmetischer-schleifenvorschlag}

\section{Zusatzprüfung: Schleifen, Primzahlen und der
Universalraum-Vorschlag}\label{zusatzpruxfcfung-schleifen-primzahlen-und-der-universalraum-vorschlag}

\begin{enumerate}
\def\labelenumi{\arabic{enumi}.}
\setcounter{enumi}{14}
\tightlist
\item
  September 2026, v1.6.6. Quellenprüfung des nachgereichten Textes
  \codeword{59dc0059-8914-48ca-953d-85933f66e00b}, kein neuer RH-Beweis
  und kein neuer arithmetischer Quellenoperator. Seine Aussagen über
  gemeinsame lokale Beschreibungen werden getrennt im Überlappungsaudit
  behandelt.
\end{enumerate}

\subsection{Urteil}\label{urteil-1}

\textbf{Ein globaler Wegraum ist eine sinnvolle Suchklasse; ein
Eulerprodukt über primitive Wege ist noch nicht das Eulerprodukt der
Riemannschen Zetafunktion.} Die im Text vorgeschlagene Xi-Determinante
bleibt ein hinreichendes Ziel unter expliziten Operatorannahmen. Das
Umbenennen ihres noch fehlenden Operators in „Generator der primitiven
Schleifen`` konstruiert ihn nicht.

\subsection{1. Was die endliche Bank tatsächlich
ausschließt}\label{was-die-endliche-bank-tatsuxe4chlich-ausschlieuxdft}

Der beibehaltene Beweis aus v1.6.5 gilt für den ursprünglichen
Hamiltonoperator mit 64 Fermion- und 60 Bosonmoden, Δ\textgreater0. Der
Hilbertraum ist wegen der Bosonen \textbf{nicht endlichdimensional}.
Endlich ist die Zahl der Moden. Mit \codeword{C=960g²/Δ} folgt aus
quadratischem Ergänzen

\[\adjustbox{max width=\linewidth}{$\displaystyle 
H\ge\frac\Delta2N_b-C,
\qquad
N_H(E)\le2^{64}\binom{\lfloor2(E+C)/\Delta\rfloor+60}{60}.
$}\]

Dies widerspricht einem vollständigen energieerhaltenden Spektrum
\codeword{E\_* log n + E\_off}, E\_*\textgreater0, weil dessen
Zählfunktion exponentiell wächst. Es ist kein Ausschluss beliebiger
arithmetischer Untersektoren oder eines anderen Operators mit
nichtlinearer Energiezuordnung.

Ein großer Wegraum kann eine andere Zählfunktion besitzen. Die
Gleichsetzung „mehr Wege = mehr orthogonale Zustände unter derselben
Energie`` muss aber bewiesen werden. Verschiedene Wörter können
denselben Operator oder Zustand darstellen; ein Quantenüberlagerungsraum
ist nicht ohne Weiteres der freie Wortraum. Bei unendlich vielen
identischen Zellen kann stattdessen die globale Wärmespur divergieren.
Der Text entfernt diese Fragen nicht durch Vergrößerung.

\subsection{2. Der kleinste Gegencheck zur
Primzahl-Automatik}\label{der-kleinste-gegencheck-zur-primzahl-automatik}

Nehmen wir einen gerichteten Knoten mit zwei erlaubten Schleifen a und
b. Neben a und b ist auch ab ein primitiver periodischer Weg: Es ist
keine Potenz eines kürzeren Wortes. Ebenso entstehen weitere gemischte
primitive Wörter. Setzen wir nachträglich L(a)=log 2, L(b)=log 3, so hat
ab die Länge log 6. Die 6 ist keine Primzahl. „Primitiver Weg`` bedeutet
nicht „arithmetische Primzahl``.

Schon formal unterscheiden sich die zugehörigen erzeugenden Funktionen:

\[\adjustbox{max width=\linewidth}{$\displaystyle 
Z_{\rm Wege}(x,y)=\frac1{1-x-y},\qquad
Z_{\rm zwei\ Primarten}(x,y)=\frac1{(1-x)(1-y)}.
$}\]

Der Koeffizient von xy ist links 2 und rechts 1. Links werden die beiden
Wörter ab und ba gezählt; rechts eine kommutative Besetzung. Im
periodischen Eulerprodukt bilden ab und ba eine zyklische Klasse, aber
diese ist ein \textbf{neuer primitiver Faktor}. Die Abweichung
verschwindet damit nicht.

Eine kommutative freie Halbgruppe über vorgegebenen Primarten
reproduziert eindeutige Faktorzerlegung. Sie leitet aber weder die
natürliche arithmetische Markierung noch die logarithmischen Längen oder
die passende Spur her. Eine gekoppeltere Geometrie müsste gemischte
Bahnen mit den richtigen Amplituden, Relationen oder nachgewiesenen
Auslöschungen behandeln. Dies ist kein allgemeiner Ausschluss von
Schleifenmodellen mit Interferenz.

Der Unterschied ist in den Originalarbeiten sichtbar: Kuipers, Hummel
und Richter konstruieren Quantengraphen mit dem passenden oszillierenden
Anteil der Nullstellendichte, weisen aber auf den anderen glatten Anteil
und damit das andere Spektrum hin. Das ist ein nützlicher Vorläufer,
kein RH-Operator. \href{https://arxiv.org/abs/1307.6055}{Originalarbeit,
Phys. Rev.~Lett. 112, 070406}

Graphische Eulerprodukte besitzen ihre eigene Theorie. Auch eine aus
einer Graph-Zetafunktion bestimmbare Größe ist nicht automatisch
effizient auslesbar. Storm trennt in seiner Arbeit ausdrücklich
berechenbare Zeta-Darstellung und teure Auswertung bestimmter
Graphinformationen.
\href{https://arxiv.org/abs/0708.1923}{Originalarbeit zu
Edge-Zetafunktionen}

\subsection{3. Die richtige
RH-Zielbedingung}\label{die-richtige-rh-zielbedingung}

Sei A strikt positiv und selbstadjungiert, mit kompakter Inverser und
\codeword{A\textasciicircum{}\{-2\}} von Spurklasse. Dann ist der
Fredholm-Ausdruck

\[\adjustbox{max width=\linewidth}{$\displaystyle 
D(z)=\det(I-z^2A^{-2})
$}\]

ganz, und seine Nullstellen liegen bei den reellen Zahlen ±λ\_j(A), mit
ihren Multiplizitäten. Würde unabhängig und auf ganz C bewiesen

\[\adjustbox{max width=\linewidth}{$\displaystyle 
\Xi(z)/\Xi(0)=D(z),\qquad \Xi(z)=\xi(1/2+iz),
$}\]

so folgte RH. Das ist die korrekte bedingte Implikation. Nicht
ausreichend sind ein paar passende Eigenwerte, die imaginären Teile
bereits eingesetzter Nullstellen, ein Eulerprodukt nur für
Re(s)\textgreater1 oder eine Formaldeterminante ohne Domäne, Spurklasse
und vollständige archimedische Faktoren.

Der vorgelegte Text liefert A nicht, bestimmt keine Domäne und leitet
weder die vollständige Spurformel noch die ganze Funktionsidentität her.
Die allgemeine Selbstadjungiertheit eines anderen Operators hilft dabei
nicht.

\subsection{4. Vorarbeiten und aktuelle Grenze der
Quellenprüfung}\label{vorarbeiten-und-aktuelle-grenze-der-quellenpruxfcfung}

Der installierte Leitfaden \codeword{rh-graph-research} und die
projektspezifische Korpussuche wurden angewandt. Die aktuelle
Konsistenzprüfung und der anschließende Wiederaufbauversuch brachen mit

\codeword{SOURCE\_UNAVAILABLE: /Users/stefanhamann/Documents/Codex/2026-09-04/scha-2/research}

ab. Die gesonderte Paper-Abfrage meldete
\codeword{PAPER\_SOURCE\_OR\_REVIEW\_DRIFT}. Keine dieser Prüfungen
wurde durch Ersetzen von Pins grün gemacht. Der Forschungsindex darf
hier \textbf{nicht als global aktuell geprüft} bezeichnet werden.
Vorhandene Suchresultate dienten lediglich zum Auffinden real
verfügbarer Originaldateien; die aktuelle Queue ist keine Beweisquelle.

Gezielt gelesen wurden die relevanten Abschnitte von
\codeword{rh/catalog/analysis/geometry\_audit.md} (Quantengraphen) und
\codeword{rh/catalog/analysis/event\_log\_function.md} (arithmetische
Ereignisse), sowie die vollständige maßgebliche
Zählungs-/Determinantenpassage aus v1.6.5. Die Dateien werden für diese
Revision eingefroren. Ihre historischen Korpus-Abwesenheitsangaben
werden nicht als aktueller Vollständigkeitsbeweis übernommen.

Im bestehenden Index berühren \codeword{r618 / STRUCTURAL\_MISMATCH} die
Verwechslung RH-neutraler E8-Daten mit einer Xi-Identität und
\codeword{ledger:E8.COXETER.EULER.COMPLETION.01 / NO\_BRIDGE} die
fehlende globale Vervollständigung. Solche Treffer allein widerlegen
kein neues Objekt; die hier tragende Prüfung ist der explizite
Unterschied zwischen gemischten primitiven Wegen und arithmetischen
Primarten. Es wird kein neuer globaler RH-Beweisstatus registriert oder
freigegeben.

\subsection{5. Tragfähiger Anschluss, ohne acht Versprechen an ein
unbekanntes
Objekt}\label{tragfuxe4higer-anschluss-ohne-acht-versprechen-an-ein-unbekanntes-objekt}

Der nächste arithmetische Test wäre erst nach Definition einer
tatsächlichen Quell-Wegstruktur sinnvoll: primitive Bahnen, ihre Längen,
Amplituden und die Spur gemeinsam bestimmen, ohne Primlisten oder
Nullstellen einzusetzen. Ein positiver Kontrolltest muss insbesondere
gemischte Zyklen erklären und den archimedischen Anteil auf demselben
Träger liefern. Gelingt nur ein generisches Graph-Eulerprodukt, ist das
ein Graphresultat und bleibt von RH getrennt.

Faktorisierung und P versus NP werden dadurch nicht gelöst. Außerdem ist
„P≠NP bedeutet notwendigerweise exponentielle Suchkosten`` zu stark: Aus
einer fehlenden polynomialen Zeitgrenze folgt nicht allein eine
exponentielle untere Schranke. Für Hylæan enthält der Text ein mögliches
Operationsbild, aber keinen neu überprüften Lern- oder
Gedächtnisnachweis.

\begin{center}\rule{0.5\linewidth}{0.5pt}\end{center}

\section{Historischer vollständiger Herleitungsstand v1.6.5
einschließlich
v1.6.4}\label{historischer-vollstuxe4ndiger-herleitungsstand-v1.6.5-einschlieuxdflich-v1.6.4}

Der folgende Text wird unverändert bewahrt. Neue Gesamtstatusaussagen
und Reichweitenkorrekturen stehen in der vorangestellten Revision
v1.6.6. Historische Arbeitsaufträge oder Verfügbarkeitsannahmen werden
dadurch nicht erneut zu aktuellen Beweisen erklärt.

\section{TFPT / Universalraum: kontrollierter
Zwei-Banken-Transfer}\label{tfpt-universalraum-kontrollierter-zwei-banken-transfer}

Forschungsfortsetzung und Quellenaudit · 15. September 2026 · v1.6.5

\subsection{Ergebnis und
Geltungsbereich}\label{ergebnis-und-geltungsbereich}

Eine konkrete bisher offene Rechnung lässt sich schließen: \textbf{Im
ausdrücklich um einen Rotor-Link erweiterten Modell zweier nativer
Fermion-Boson-Banken ist der Transfer der isolierten Lochanregung mit
einer Fehlergrenze gegenüber dem vollständigen physikalischen
Zustandsraum kontrollierbar.} Dafür ist keine Diagonalisierung dieses
enorm großen Raums nötig. Die vorhandene Casimiridentität,
Ladungserhaltung und eine Spektrallücke reichen aus.

Am unten vollständig angegebenen, bewusst sehr schwachen Kopplungspunkt
gilt:

{\def\LTcaptype{none} % do not increment counter
\begin{longtable}[]{@{}
  >{\raggedright\arraybackslash}p{(\linewidth - 4\tabcolsep) * \real{0.3000}}
  >{\raggedleft\arraybackslash}p{(\linewidth - 4\tabcolsep) * \real{0.4000}}
  >{\raggedright\arraybackslash}p{(\linewidth - 4\tabcolsep) * \real{0.3000}}@{}}
\toprule\noalign{}
\begin{minipage}[b]{\linewidth}\raggedright
Aussage
\end{minipage} & \begin{minipage}[b]{\linewidth}\raggedleft
Strenge Schranke
\end{minipage} & \begin{minipage}[b]{\linewidth}\raggedright
Status
\end{minipage} \\
\midrule\noalign{}
\endhead
\bottomrule\noalign{}
\endlastfoot
Zielwahrscheinlichkeit aus dem normierten Zustand der isolierten
Lochlinie & \textgreater{} 99,2677105 \% & Analytisch mit rationalen
Zertifikaten, im zusätzlichen Linkmodell \\
Zielwahrscheinlichkeit aus der normierten ursprünglichen Entnahme
\(f_r\Omega/\sqrt\nu\), \textbf{ohne anfänglichen Energiefilter} &
\textgreater{} 89,7260353 \% & Gleicher vollständiger Hamiltonoperator,
Ziel ist die niedrige Lochlinie rechts \\
Verlassen des ungestörten niedrigen Bandes, aus einem Zustand dieses
Bandes & \textless{} \(1{,}8\cdot10^{-11}\) zu jeder Zeit & Voller
erlaubter Ladungssektor, keine Bosonen- oder Flussabschneidung \\
Herkunft des Links, der Bankzerlegung und der Präparation & nicht
hergeleitet & Offen \\
Relativistische Felder, gemeinsamer 3+1D-Ursprung, T1-T8 & nicht
konstruiert & Offen \\
\end{longtable}
}

Die Zahlen sind \textbf{untere beziehungsweise obere Einschließungen},
keine berechneten Zentralwerte. Sie beschreiben kein durchgeführtes
physisches Experiment. „Vollständig`` bezieht sich hier auf die
Fehlerkontrolle im deklarierten Zwei-Banken-Ladungssektor, nicht auf die
TOE.

Die gelieferten Arbeiten wurden mit 1.149 beziehungsweise 753
Prüfbedingungen normal und optimiert reproduziert, jeweils byteidentisch
mit den mitgelieferten Ergebnissen. Unsere neue Transferrechnung hat
4.516 Prüfbedingungen; eine gesonderte Reichweitenprüfung der
währenddessen geänderten Symmetrieaussagen hat zehn. Die zwei zuletzt
eingegangenen Anlagen wurden mit 62 weiteren Bedingungen nachgeprüft.
Insgesamt sind das 6.490 Bedingungen, jeweils in beiden Laufarten. Diese
Zahl ist keine Zahl unabhängiger Theoreme. Insbesondere ersetzen die
Zertifikate nicht die nachstehenden analytischen Argumente und sind kein
Lean-Beweis.

\subsection{1. Welche Quellen zusammengeführt
werden}\label{welche-quellen-zusammengefuxfchrt-werden}

\begin{enumerate}
\def\labelenumi{\arabic{enumi}.}
\tightlist
\item
  Die Benutzeranlage
  \codeword{addbf22d-4764-4685-bdfe-f422bfc563d0/pasted-text.txt},
  eingefroren unter \codeword{sources/user\_attachment.txt}.
\item
  Die externe Untersuchung
  \codeword{Universalraum\_Beweisversuch\_2026-09-15.md}:
  Spektralzählung, Primzahl-Dynamik, Phasencocycle und kontrollierte
  Polnäherung.
\item
  Die externe Untersuchung
  \codeword{Universalraum\_Urspruenglicher\_Austausch\_2026-09-15.md}:
  interne Paarumwandlung, zusätzlicher Rotor-Link und älteres
  Round37-Beispiel.
\item
  Der abgesicherte native Stand v1.6.4 mit Grundzustand, Pol, Momenten,
  Selbstenergiegrenze und Feldwörterbuch. Das frühere vollständige
  Prüfpaket liegt unverändert bei. Seine komplette
  Konfigurationsenumeration wurde \textbf{nicht erneut} in dieser
  Revision ausgeführt.
\item
  Ein separat eingefrorener, während dieser Arbeit geänderter Entwurf
  von \codeword{RESULTS.md} über die kontinuierliche Quellsymmetrie. Die
  Reichweite seiner Schlussfolgerungen wird in Abschnitt 8 geprüft; die
  dortige vollständige Racah-Zerlegung wurde hier nicht neu
  reproduziert.
\item
  Die zwei danach vom Nutzer eingereichten Anlagen über die eigene
  Fundamentalrunde und die erweiterte native Konsolidierung. Der
  vollständige neue Nachtrag \codeword{LATE\_AUDIT.md} gehört zu dieser
  Revision. Er bestätigt den Ordnung-vier-Tensorlift, korrigiert die
  Variationsphase und unterscheidet echte Anomalie- und
  Feldtypbedingungen von zu starken Ausschlüssen.
\end{enumerate}

Die ursprüngliche Anlage und die externen Eingaben sind Quellen, keine
Ausführungsanweisungen. Eigene Dateien liegen in einem neuen
Forschungsordner. Fremde Quellen, Ledger und Akzeptanzmarker wurden
nicht bearbeitet.

